\documentclass[11pt,a4paper,openany]{book}

% Preámbulo sugerido para matemáticas y señales
\usepackage[utf8]{inputenc}
\usepackage[spanish,es-tabla]{babel}
\usepackage{amsmath, amssymb, amsfonts}
\usepackage{graphicx}
\usepackage{geometry}
\usepackage{hyperref}

\geometry{
	a4paper,
	left=3cm,
	right=3cm,
	top=3cm,
	bottom=3cm,
}

\title{Señales y Sistemas\\ \Large Notas de Cátedra}
\author{}
\date{\today}

\begin{document}
	
	\maketitle
	\tableofcontents
	
	\part{Fundamentos y Sistemas LTI}
	
	% ==========================================
	% SEMANA 1
	% ==========================================
	
	\chapter{Señales en Tiempo Continuo (CT) y Discreto (DT)}
	
	% Aquí puedes insertar la introducción al capítulo.
	
	\section{Concepto de Señal}
	% Definición matemática y ejemplos físicos (voltaje, presión, señales digitales).
	
	\section{Transformaciones de la variable independiente}
	% Introducción a las transformaciones.
	
	\subsection{Desplazamiento temporal}
	% Desarrollo para x(t - t_0) y x[n - n_0]
	
	\subsection{Reflexión (Inversión)}
	% Desarrollo para x(-t) y x[-n]
	
	\subsection{Escalamiento temporal}
	% Desarrollo para x(at) (compresión y expansión)
	
	\section{Señales Periódicas}
	% Definición de periodo fundamental T y N.
	
	\section{Señales Elementales}
	% Introducción a las señales básicas.
	
	\subsection{Exponenciales complejas y senoidales}
	% Diferencias de periodicidad entre CT y DT.
	
	\subsection{La función impulso unitario y escalón unitario}
	% Desarrollo sobre \delta y u. Relación entre ambas (derivada/integral y diferencia/suma).
	
	
	\chapter{Definición y Propiedades de los Sistemas}
	
	% Aquí puedes insertar la introducción al capítulo.
	
	\section{Interconexiones de sistemas}
	% Serie (cascada), paralelo y realimentación.
	
	\section{Propiedades Fundamentales}
	% Definiciones y métodos de verificación.
	
	\subsection{Memoria}
	% Sistemas instantáneos vs. sistemas con memoria.
	
	\subsection{Invertibilidad}
	% Concepto de sistema inverso.
	
	\subsection{Causalidad}
	% Dependencia de valores presentes y pasados.
	
	\subsection{Estabilidad (BIBO)}
	% Entradas acotadas producen salidas acotadas.
	
	\subsection{Invarianza en el tiempo}
	% El comportamiento no cambia con el tiempo.
	
	\subsection{Linealidad}
	% Principios de aditividad y homogeneidad (superposición).
	
	
	% ==========================================
	% SEMANA 2
	% ==========================================
	
	\chapter{La Suma de Convolución (Tiempo Discreto)}
	
	% Aquí puedes insertar la introducción al capítulo.
	
	\section{Representación de señales DT en términos de impulsos}
	% Cómo "desarmar" x[n] usando \delta[n-k].
	
	\section{Respuesta al Impulso ($h[n]$)}
	% Definición como la huella dactilar del sistema.
	
	\section{La Suma de Convolución}
	% Desarrollo conceptual.
	\begin{equation}
		y[n] = \sum_{k=-\infty}^{\infty} x[k]h[n-k]
	\end{equation}
	
	\section{Cálculo práctico}
	% Método gráfico (reflejar, desplazar y multiplicar) y método analítico.
	
	
	\chapter{La Integral de Convolución y Propiedades de LTI}
	
	% Aquí puedes insertar la introducción al capítulo.
	
	\section{Representación de señales CT en términos de impulsos}
	% El concepto del límite de pulsos estrechos.
	
	\section{La Integral de Convolución}
	% Desarrollo conceptual.
	\begin{equation}
		y(t) = \int_{-\infty}^{\infty} x(\tau)h(t-\tau) d\tau
	\end{equation}
	
	\section{Propiedades de los sistemas LTI}
	% Introducción a las propiedades.
	
	\subsection{Conmutatividad, Distributividad y Asociatividad}
	% Desarrollo de las propiedades.
	
	\subsection{Caracterización de propiedades LTI a través de la respuesta al impulso}
	% ¿Cómo saber si un sistema LTI es causal o estable mirando solo su respuesta al impulso h(t) o h[n]?
	
	\section{Respuesta al Escalón}
	% Relación con la respuesta al impulso.
	
	% ==========================================
	% BIBLIOGRAFÍA Y EJERCICIOS
	% ==========================================
	
	\cleardoublepage
	\phantomsection
	\addcontentsline{toc}{chapter}{Bibliografía y Ejercicios Recomendados}
	\chapter*{Bibliografía y Ejercicios Recomendados}
	
	\textbf{Bibliografía específica para este bloque:}
	\begin{itemize}
		\item \textbf{Oppenheim:} Capítulo 1 (Secciones 1.1 a 1.6) y Capítulo 2 (Secciones 2.1 a 2.4).
	\end{itemize}
	
	\textbf{Ejercicios recomendados:}
	\begin{itemize}
		\item Problemas de fin de capítulo 1.21 a 1.27 (propiedades de sistemas).
		\item Problemas de fin de capítulo 2.21 a 2.23 (convoluciones básicas).
	\end{itemize}
	
% ==========================================
% INICIO DEL BLOQUE II
% ==========================================

\part{Análisis de Fourier}

% ==========================================
% SEMANA 3
% ==========================================

\chapter{Series de Fourier (Señales Periódicas)}

\section{Serie de Fourier en Tiempo Continuo (CTFS)}
% Desarrollo de los siguientes temas:
% - Respuesta de sistemas LTI a exponenciales complejas: Autofunciones y autovalores H(j\omega).
% - Representación de señales periódicas: Ecuaciones de síntesis y análisis.
% - Determinación de coeficientes (a_k): Ortogonalidad de las exponenciales complejas.
% - Fenómeno de Gibbs: Comportamiento en discontinuidades.
% - Propiedades: Linealidad, desplazamiento, conjugación y simetría.
% - Relación de Parseval.

\section{Serie de Fourier en Tiempo Discreto (DTFS)}
% Desarrollo de los siguientes temas:
% - Diferencia fundamental con CTFS (periodicidad de las exponenciales discretas, N coeficientes).
% - Ecuaciones de síntesis y análisis (sumatorias finitas).
% - Propiedades de la DTFS.
% - Sistemas LTI y Series de Fourier: Filtrado de señales periódicas.


% ==========================================
% SEMANA 4
% ==========================================

\chapter{Transformada de Fourier en Tiempo Continuo (CTFT)}

\section{Representación de Señales Aperiódicas}
% Extensión de la Serie de Fourier: El límite cuando el periodo T \to \infty.

\subsection{El Par de la Transformada de Fourier}
% Desarrollo conceptual y condiciones de convergencia (Dirichlet).

\textbf{Ecuación de Análisis:}
\begin{equation}
	X(j\omega) = \int_{-\infty}^{\infty} x(t) e^{-j\omega t} dt
\end{equation}

\textbf{Ecuación de Síntesis:}
\begin{equation}
	x(t) = \frac{1}{2\pi} \int_{-\infty}^{\infty} X(j\omega) e^{j\omega t} d\omega
\end{equation}

% Ejemplos clásicos: Pulso rectangular y función Sinc.

\section{Propiedades y Sistemas LTI en Frecuencia}
% Desarrollo de propiedades clave (desplazamiento, escalamiento, diferenciación/integración).
% Propiedad fundamental de Convolución: y(t) = x(t) * h(t) \iff Y(j\omega) = X(j\omega)H(j\omega).
% Sistemas descritos por ecuaciones diferenciales (obtención de H(j\omega)).


% ==========================================
% SEMANA 5
% ==========================================

\chapter{Transformada de Fourier en Tiempo Discreto (DTFT)}

\section{Definición y Pares Comunes}
% Definición y justificación de por qué X(e^{j\omega}) es siempre periódica (periodo 2\pi).
% Pares comunes: Impulso, escalón, señales exponenciales causales.
% Dualidad entre las transformaciones (CTFS, DTFS, CTFT, DTFT).

\section{Propiedades de la DTFT y Repaso Integral}
% Convolución en tiempo discreto.
% Sistemas descritos por ecuaciones de diferencias (obtención de H(e^{j\omega})).
% Matriz comparativa de propiedades y resolución de dudas.


% ==========================================
% SEMANA 6
% ==========================================

\chapter{Evaluación del Bloque I y II}

\section{Taller de Problemas}
% Resolución de exámenes anteriores. 
% Énfasis en linealidad/causalidad, convoluciones gráficas y espectros de magnitud/fase.

\section{Primer Examen Parcial}
% Temario: Capítulos 1 al 5 (Oppenheim).
% Formato: 2 problemas LTI/convolución, 2 problemas Fourier (CT y DT).


% ==========================================
% BIBLIOGRAFÍA Y EJERCICIOS - BLOQUE II
% ==========================================

\cleardoublepage
\phantomsection
\addcontentsline{toc}{chapter}{Bibliografía y Ejercicios Recomendados - Bloque II}
\chapter*{Bibliografía y Ejercicios Recomendados - Bloque II}

\textbf{Bibliografía específica para este bloque:}
\begin{itemize}
	\item \textbf{Oppenheim, Willsky y Nawab:} Capítulos 3, 4 y 5 completos.
\end{itemize}

\textbf{Ejercicios clave recomendados:}
\begin{itemize}
	\item 3.22 (Cálculo de coeficientes de serie).
	\item 4.21 (Pares de transformada en tiempo continuo).
	\item 5.21 (Pares de transformada en tiempo discreto).
\end{itemize}

% ==========================================
% INICIO DEL BLOQUE III
% ==========================================

\part{Caracterización y Muestreo}

% ==========================================
% SEMANA 7
% ==========================================

\chapter{Respuesta en Magnitud y Fase}

\section{Análisis de la respuesta en frecuencia}
% Representación de H(j\omega) y H(e^{j\omega}) en forma polar.

\section{Efectos de la Fase}
% Concepto de retraso de grupo (group delay) y distorsión de fase. Importancia de la fase lineal.

\section{Escalas Logarítmicas}
% Introducción a los diagramas de Bode (magnitud en dB y fase en grados).

\section{Relación Tiempo-Frecuencia}
% Cómo el ancho de banda en frecuencia afecta la rapidez de respuesta (tiempo de subida) en el dominio del tiempo.


\chapter{Filtros y Sistemas de Segundo Orden}

\section{Filtros Ideales}
% Pasabajos, pasaaltos y pasabanda. La imposibilidad física de los filtros ideales (causalidad).

\section{Filtros No Ideales}
% Concepto de banda de paso, banda de transición y banda de rechazo.

\section{Sistemas de primer y segundo orden}

\subsection{Tiempo continuo}
% Circuitos RC, RLC y su respuesta en frecuencia.

\subsection{Tiempo discreto}
% Ecuaciones de diferencias de primer y segundo orden.

\section{Análisis de resonancia y amortiguamiento}
% Desarrollo conceptual.


% ==========================================
% SEMANA 8
% ==========================================

\chapter{Muestreo en el Dominio del Tiempo}

\section{Representación del muestreo como una multiplicación}
% El tren de impulsos unitarios como señal de muestreo.

\section{Efecto en el espectro}
% Replicación del espectro de la señal original en múltiplos de la frecuencia de muestreo \omega_s.

\section{El Teorema de Nyquist-Shannon}
% Condición para evitar el solapamiento (aliasing): \omega_s > 2\omega_M.

\section{Aliasing}
% Qué ocurre cuando se submuestrea una señal (ejemplos visuales y auditivos).


\chapter{Reconstrucción y Procesamiento Digital}

\section{Reconstrucción de señales}
% El filtro pasabajos ideal como interpolador.

\section{Interpolación práctica}
% Retención de orden cero (Zero-Order Hold) y de primer orden.

\section{Procesamiento de tiempo continuo en tiempo discreto}

\subsection{Diagrama de bloques}
% C/D (Conversor Continuo a Discreto) \to Sistema LTI Discreto \to D/C (Conversor Discreto a Continuo).

\subsection{Diseño de un sistema digital}
% Cómo diseñar un sistema digital que emule un filtro analógico.

\section{Diezmación e Interpolación}
% Conceptos básicos de cambio de la frecuencia de muestreo en el dominio discreto.


% ==========================================
% BIBLIOGRAFÍA Y EJERCICIOS - BLOQUE III
% ==========================================

\cleardoublepage
\phantomsection
\addcontentsline{toc}{chapter}{Bibliografía y Ejercicios Recomendados - Bloque III}
\chapter*{Bibliografía y Ejercicios Recomendados - Bloque III}

\textbf{Bibliografía específica para este bloque:}
\begin{itemize}
	\item \textbf{Oppenheim, Willsky y Nawab:} Capítulo 6 (Secciones 6.1 a 6.5) y Capítulo 7 (Secciones 7.1 a 7.5).
\end{itemize}

\textbf{Ejercicios clave recomendados:}
\begin{itemize}
	\item 6.21 (Diagramas de magnitud y fase).
	\item 7.21 y 7.22 (Cálculo de frecuencias de Nyquist y reconstrucción).
\end{itemize}

% ==========================================
% INICIO DEL BLOQUE IV
% ==========================================

\part{Transformadas Complejas ($s$ y $z$)}

% ==========================================
% SEMANA 9
% ==========================================

\chapter{Definición y Región de Convergencia (ROC)}

\section{De Fourier a Laplace}
% Generalización de la transformada mediante la variable compleja s = \sigma + j\omega.

\section{Definición de la Transformada bilateral}
% Desarrollo conceptual.
\begin{equation}
	X(s) = \int_{-\infty}^{\infty} x(t) e^{-st} dt
\end{equation}

\section{La Región de Convergencia (ROC)}
% Por qué el álgebra no es suficiente; importancia de la ROC para identificar la señal original.

\section{Propiedades de la ROC}
% Relación con la causalidad y la duración de la señal (señales de lado derecho, izquierdo y biltaterales).

\section{Diagrama de Polos y Ceros}
% Visualización de la función de transferencia H(s) en el plano s.

\chapter{Propiedades y Estabilidad de Sistemas LTI}

\section{Propiedades de Laplace}
% Linealidad, desplazamiento, escalamiento y el teorema del valor inicial/final.

\section{Transformada de Laplace Inversa}
% Uso de fracciones parciales para volver al dominio del tiempo.

\section{Análisis de Sistemas LTI}

\subsection{Función de transferencia $H(s)$}
% Desarrollo de la función de transferencia.

\subsection{Causalidad}
% La ROC es un semiplano a la derecha del polo más a la derecha.

\subsection{Estabilidad}
% Un sistema LTI es estable si y solo si la ROC incluye el eje j\omega (equivalente a decir que todos los polos están en el semiplano izquierdo para sistemas causales).


% ==========================================
% SEMANA 10
% ==========================================

\chapter{Definición y el Plano Z}

\section{Generalización en tiempo discreto}
% La variable compleja z = re^{j\omega}.

\section{Definición de la Transformada Z}
% Desarrollo conceptual.
\begin{equation}
	X(z) = \sum_{n=-\infty}^{\infty} x[n] z^{-n}
\end{equation}

\section{Relación entre el plano $s$ y el plano $z$}
% El mapeo del eje imaginario de s al círculo unitario de z.

\section{ROC en el plano Z}
% Representación como anillos concéntricos. Propiedades de la ROC para secuencias finitas, derechas e izquierdas.

\chapter{Propiedades, Inversa y Estabilidad}

\section{Propiedades de la Transformada Z}
% Desplazamiento temporal (el operador z^{-1} como delay), convolución y diferenciación en el dominio z.

\section{Transformada Z Inversa}
% Método de expansión en fracciones parciales y división larga (series de potencias).

\section{Caracterización de Sistemas LTI en el dominio $z$}

\subsection{Estabilidad y Causalidad}
% Un sistema LTI discreto es estable si la ROC incluye el círculo unitario (|z|=1).
% Para sistemas causales, esto implica que todos los polos deben estar dentro del círculo unitario.


% ==========================================
% BIBLIOGRAFÍA Y EJERCICIOS - BLOQUE IV
% ==========================================

\cleardoublepage
\phantomsection
\addcontentsline{toc}{chapter}{Bibliografía y Ejercicios Recomendados - Bloque IV}
\chapter*{Bibliografía y Ejercicios Recomendados - Bloque IV}

\textbf{Bibliografía específica para este bloque:}
\begin{itemize}
	\item \textbf{Oppenheim, Willsky y Nawab:} Capítulo 9 (Laplace) y Capítulo 10 (Z-Transform).
\end{itemize}

\textbf{Ejercicios clave recomendados:}
\begin{itemize}
	\item 9.21 y 9.22 (Identificar señales a partir de polos y ROC).
	\item 10.21 y 10.22 (Transformada Z y diagramas de polos y ceros).
\end{itemize}

% ==========================================
% INICIO DEL BLOQUE V
% ==========================================

\part{Aplicaciones, Filtros y Simulación}

% ==========================================
% SEMANA 11
% ==========================================

\chapter{Muestreo en el Dominio de la Frecuencia}

\section{La Transformada Discreta de Fourier (DFT)}
% Cómo representar el espectro de una señal mediante un número finito de muestras.

\section{Relación entre la DTFT y la DFT}
% Muestreo de la DTFT en el círculo unitario.

\section{Resolución espectral}
% El efecto del enventanado (windowing) y la longitud del registro de datos.

\chapter{Algoritmos y Procesamiento Rápido}

\section{Introducción a la FFT}
% Concepto de diezmación en tiempo y frecuencia (eficiencia computacional).

\section{Convolución Circular vs. Lineal}
% Por qué la multiplicación en el dominio de la DFT no es exactamente igual a la convolución lineal y cómo corregirlo (zero-padding).


% ==========================================
% SEMANA 12
% ==========================================

\chapter{Estructuras de Filtros Digitales}

\section{Representación mediante diagramas de bloques}
% Sumadores, ganancias y retrasos ($z^{-1}$).

\section{Formas Directas}
% Forma Directa I y Forma Directa II (Canónica).

\section{Efectos de la cuantización}
% Breve introducción a los errores por precisión finita.

\chapter{Filtros FIR e IIR}

\section{Filtros de Respuesta al Impulso Finita (FIR)}
% Propiedad de fase lineal y diseño por el método de ventanas (Hamming, Hann, Blackman).

\section{Filtros de Respuesta al Impulso Infinita (IIR)}
% Diseño a partir de prototipos analógicos (Butterworth, Chebyshev) mediante la Transformación Bilineal.


% ==========================================
% SEMANA 13
% ==========================================

\chapter{Representación en Variables de Estado (State-Space)}

\section{Concepto de Estado}
% Cómo pasar de ecuaciones diferenciales/diferencias de orden N a un sistema de N ecuaciones de primer orden.

\section{Ecuaciones de estado y de salida}
% Matrices A, B, C, D.

\section{Ventajas de la representación de estado}
% Modelado multivariable y simulación numérica eficiente.

\chapter{Simulación Computacional}

\section{Algoritmos de integración numérica}
% Euler y Runge-Kutta (conceptos básicos aplicados a señales).

\section{Simulación de sistemas realimentados}
% Estabilidad numérica y retardo de transporte.

\section{Uso de herramientas de software}
% Breve repaso de cómo implementar estos modelos en entornos tipo MATLAB o Python (SciPy).


% ==========================================
% SEMANA 14
% ==========================================

\chapter{Taller de Integración}

\section{Resolución de un problema ``extremo-a-extremo''}
% Desde la toma de una señal continua, su muestreo, filtrado digital y reconstrucción.

\section{Análisis de compromiso (\textit{trade-offs})}
% Velocidad de muestreo vs. carga computacional.

\chapter{Sesión de Consultas}

\section{Repaso General}
% Repaso de los conceptos más complejos del Bloque IV (Z y Laplace) y Bloque V.

\section{Simulacro de examen parcial}
% Preparación para la evaluación.


% ==========================================
%
\end{document}